As we describe in \secref{three_methods}, the IJ covariance is particularly
well-motivated in settings when the MAP $\thetahat$ is either unavailable or
unreliable, and one must turn to MCMC to integrate out a large number of
nuisance parameters.  In this section we present a simple simulated example of
such a case.
% A similar example on a real dataset can be found in PUT IN PILOTS
% can be found in REF APPENDIX.

Consider the following simple random effects model with responses. For the data,
we assume we have responses $y_n \in \rdom{}$, regressors $r_n \in \rdom{}$, and
group indicator vectors $z_n \in \{0, 1\}^K$, for $n=1,\ldots,N$ and some fixed
number of groups $K$.  We assume that all three are conditionally IID, so $\x_n
= (y_n, r_n, z_n)$, and draw according to the model
%
\begin{align}
\varepsilon_n \indep \normdist(\cdot \vert 0, 1)
\quad\quad
r_n \indep \normdist(\cdot \vert 0, 1)
\quad\quad
y_n = r_n + r_n^2 \varepsilon_n.\eqlabel{simple_sim_dgp}
\end{align}
%
In our experiments we used $K = \simNumRe$ and $N = \simNumObs$.

Suppose we now form a misspecified model of this data. Define the regression
parameter $\beta \in \rdom{}$, random effect vector $\mu \in \rdom{K}$, and
variances $\sigma_y^2$ and $\sigma_\mu^2$.  We will take the priors on $\mu_k$
to be independent given $\sigma_\mu$, and the priors on $\beta$, $\sigma_\y$ and
$\sigma_\mu$ to be diffuse; we will shortly be using the default diffuse priors
provided by \texttt{rstanarm}. Then we can model
%
\begin{align}\eqlabel{simple_sim_model}
%
y_n | z_n, r_n, \beta, \mu, \sigma_y^2 \indep&
    \normdist(\cdot | \beta r_n + z_n^T \mu, \sigma_y^2)
    \quad\textrm{ and }\quad
\mu_k | \sigma_\mu \indep
    \normdist(\cdot | 0, \sigma_\mu^2).
%
\end{align}
%
The model in \eqref{simple_sim_model} is misspecified because it does not model the
heteroskedasticity induced by $r_n^2 \varepsilon_n$ in \eqref{simple_sim_dgp}.
Note also that the group indicators $\z_n$ have no effect on the real data
generating process, and so estimates of very small $\sigma_\mu$ should be
supported by the data --- and the log likelihood should be degenerate at the
point $\sigma_\mu$, as we describe in detail below.

How good is the MAP estimate for this model?  The answer depends on which
variables are marginalized in the definition of $\theta$.  Suppose, for example,
we form the ``global'' MAP, optimizing over all the parameters,
including the random effects:
%
\def\badmap{\textrm{glob}}
\begin{align*}
%
\hat\beta^{\badmap}, \hat\sigma_y^{\badmap},
\hat\sigma_\mu^{\badmap}, \hat\mu^{\badmap} :=
    \argmax \ell(\xvec | \beta, \sigma_y, \sigma_\mu, \mu) +
            \log \pi(\mu | \sigma_\mu) +
            \log \pi(\beta, \sigma_y, \sigma_\mu).
%
\end{align*}
%
Optimizing over all parameters including the random effects is prone to
under-estimate the variance $\sigma_y$, particularly when $K$ is large relative
to $N$.  A simplified version of this phenomenon is the known as the
Neyman-Scott paradox \citep{freedman:2005:neymanscott}.  For this reason, one
does not typically optimize over $\mu$ for such a model.  Rather, one optimizes
the marginal posterior (or an approximation thereof):
%
\begin{align}\eqlabel{singular_map}
%
\hat\beta, \hat\sigma_y, \hat\sigma_\mu :=
    \argmax & \log \int p(\xvec | \beta, \sigma_y, \sigma_\mu, \mu)
        \pi(\mu | \sigma_\mu) d\mu  +
       \log \pi(\beta, \sigma_y, \sigma_\mu)
      \nonumber\\
= \argmax & \sum_{k=1}^K \ell(\xvec_k | \beta, \sigma_y, \sigma_\mu) +
\log \pi(\beta, \sigma_y, \sigma_\mu),
%
\end{align}
%
where $\xvec_k$ denotes $\{\x_n: z_{nk} = 1\}$, the observations in group $k$.
For Gaussian models such as this one, the needed integral in
\eqref{singular_map} can be computed in closed form. In the general case, there
exist non-MCMC methods such as marginal maximum likelihood estimation (MMLE) use
the Laplace approximation to approximate the integral
\citep{bates:2015:fitting}. However, it is possible to generate data $\xvec$ so
that even the MAP of \eqref{singular_map} estimates $\hat\sigma_\mu = 0$ for
sufficiently non-informative priors
\citep{chung:2012:avoiding,eager:2017:mixedterrible}.

In this case, we have designed the model and data generating process so that the
MMLE package \texttt{lme4} results in a numerically singular fit
($\hat\sigma_\mu = 0$) approximately half the time. When $\hat\sigma_\mu = 0$,
the second derivative of the log posterior is degenerate at the optimum, and
classical methods for approximating frequentist variances cannot be directly
applied.  In a random effects model for which even MMLE methods fail, MCMC is a
natural recourse. If an estimate of the frequentist variability of the posterior
mean is required, the IJ covariance is then a natural choice.


% Further, the heteroskedasticity in the residual
% variance, together with the modeling assumption of homoskedasticity, means that
% the Bayesian and frequentist covariances are expected to be different.


% In short, if we take $g(\theta) = \theta := \beta$ in this problem, then
% \eqref{nonsingular_map} gives us reason to believe that $\thetahat \approx
% \expect{\post}{\theta}$ and that both are approximately normally distributed,
% due to the concentration of the marginal log likelihood $\sumn \ell(\x_n |
% \theta)$.  However, we cannot easily compute the objective function or its
% derivatives, since $\ell(\x_n | \theta)$ involves an intractable integral.
% Nevertheless, we can run MCMC on all the parameters and use the IJ to
% approximate the frequentist variance of the limiting distribution of
% $\expect{\post}{\theta}$.  It is in precisely such cases that the IJ covariance
% is most useful.

\SingularResultsGraph{}

We simulated a single dataset for which the \texttt{lme4} fit happened to be
singular, and computed $\simNumMCMCDraws$ MCMC draws from the model of
\eqref{simple_sim_model} using \texttt{rstanarm}.  With these samples, we
computed the Bayesian posterior covariance and the IJ covariance.  We then
simulated $\simNumSims$ different datasets, ran the same MCMC procedure for
each, computed the posterior mean for each simulation, and computed the
covariance of the posterior means over the simulations.  Ideally, we expect the
IJ covariance to provide a good approximation to the simulation-based covariance
estimate.

The results are shown in \figref{simple_sim_result}.  As shown in the left plot,
the posterior variances of $\log \sigma_y$ and of $\beta$ are dramatically
less than the frequentist variances of $\expect{\post}{\log \sigma_y}$ and
$\expect{\post}{\beta}$.  The posterior of the log random effect variance, $\log
\sigma_\mu$, is extremely non-normal, and the IJ may not expected to perform
well, theoretically.  Nevertheless, we show that the IJ covariance is correct
despite the posterior variance being extremely inflated.